\chapter*{How to read this book}
\addcontentsline{toc}{chapter}{How to read this book}
\markboth{How to read this book}{How to read this book}

\section*{What this is}

This book consolidates the entire mathematical content of \textsc{tatiana}
into a single document. Its sources are the files in \texttt{DOCS/}, the
poster sections in \texttt{POSTER/}, and the specification documents added
on the branch \texttt{fix-11-13-and-curvature-decisions}. It is
\emph{implementation-free}: no C\raisebox{0.2ex}{\small ++}, no Python.
Where the implementation determines a mathematical fact --- and it
sometimes does --- the fact is stated abstractly here and the discrepancy
is recorded in Appendix~\ref{app:discrepancy}.

\section*{Why every claim carries a tag}

The source documents \emph{contradict one another}. The internal audit
(\texttt{DOCS/AUDIT\_SCRUTINY\_AND\_BOOK.md}) records fourteen findings
against the earlier mathematical description, several of which are
outright refutations: a proposition false in the deployed case, a theorem
stated more generally than it is proved, a repair for $\Hone$ that is
half circular. A smooth merge of these sources would present refuted
material in the same confident register as proved material, and a reader
studying it would learn falsehoods.

So every substantive claim in this book carries one of six tags:

\begin{center}
\begin{tabular}{ll}
\toprule
\Proved & A complete proof is given here, from stated hypotheses.\\
\ProvedCond & Proved, but only under hypotheses the deployed system may not meet.\\
\Implemented & The system computes this; no proof of correctness exists.\\
\Proposed & A design intention. Not built, not proved.\\
\Contested & The audit disputes this. The dispute is stated in place.\\
\Refuted & Tested and failed, or shown false. Retained \emph{with} the refutation.\\
\bottomrule
\end{tabular}
\end{center}

Refuted material is kept rather than deleted. Knowing why a construction
failed is most of what you need in order not to rebuild it.

\section*{The three layers}

Each central object is presented three times over, so that a single
reading can be shallow or deep:

\begin{enumerate}
  \item a boxed \textbf{Intuition} in plain language, with no symbols;
  \item the formal \textbf{Definition}, \textbf{Construction} or
        \textbf{Proposition}, with proof;
  \item a \textbf{Worked example} on an explicit small complex, carried
        through to numbers.
\end{enumerate}

The worked examples are cumulative. A single running example --- a
seven-organ complex called $\Kc_7$, introduced in
Chapter~\ref{ch:complex} --- is reused throughout, so that the same
object is measured by every instrument the book builds. Where an
exercise appears, it is because the step is short and doing it yourself
is faster than reading it.

\section*{Provenance}

Grey boxes marked \textsf{\textbf{Source}} record which file a passage
descends from, and what was changed in merging it. They exist so that
any statement here can be traced back and checked against the original,
and so that when the sources disagree you can see \emph{that} they
disagree rather than only seeing the winner.

\section*{Reading order}

Parts I and II are the conceptual spine and everything else depends on
them; read them in order. Parts III through VII may be read in any
order once the spine is in place. Part VIII is the honest boundary and
should be read before quoting any result from this book to anyone else.
Parts IX and X are forward-looking: open problems, and the quantum
implementation programme.
